% Journal of Sports Economics submission. SAGE asks for a plain double-spaced
% manuscript (12pt Times, 1-inch margins) with APA 7th author-date references;
% apacite (with natbibapa) provides the APA citation commands while keeping the
% \citep/\citet calls already in the text.
\documentclass[12pt,USenglish]{article}

\usepackage[utf8]{inputenc}
\usepackage[T1]{fontenc}
\usepackage[margin=1in]{geometry}
\usepackage{mathptmx}            % Times for text and math
\usepackage{microtype}
\usepackage{setspace}
\usepackage{graphicx}
\usepackage{booktabs}
\usepackage{multirow}
\usepackage{array}
\usepackage{longtable}
\usepackage{amsmath}
\usepackage{amssymb}
\usepackage[hidelinks]{hyperref}
\usepackage[natbibapa]{apacite}  % APA 7th author-date; defines \citep/\citet

\doublespacing

\begin{document}

\title{The Limits of Foresight in Coaching Hires:\\
Pre-Hire Predictability of Involuntary Separation and the\\
Internal-Promotion Hazard in the NFL}

% Author identity is carried on the separate title_page.tex (blinded review).
\author{}
\date{}

\maketitle

\begin{abstract}
\noindent When an NFL franchise hires a head coach, how much of the eventual firing is already foreseeable? Prior work estimates the firing hazard from performance realized during a coach's tenure; we instead estimate the ex-ante predictability of involuntary separation from the information set available on hiring day. Across 371 head coaching stints (1970--2025), we fit a cause-specific Cox model that treats voluntary departures as a competing risk and selects covariates by Cox-native stability selection. Predictability is bounded: a Harrell concordance near 0.60 and an integrated Brier score of 0.176 leave most of a firing unforeseeable when the contract is signed. One structural signal is robust. Internal promotions face a higher firing hazard than external hires (hazard ratio 1.52), and the premium persists after conditioning on observable franchise distress. The broad on-field record that hiring searches scrutinize forecasts dismissal no better than chance.
\end{abstract}

\medskip
\noindent\textbf{Keywords:} head coaching labor market, involuntary separation, survival analysis, competing risks, internal promotion

\newpage

\section{Introduction}

Hiring a head coach is one of the most consequential decisions an NFL franchise makes, and history shows how often it is undone. Of the 371 stints we study in the modern era (1970--2025), 84\% ended in an involuntary parting, and the median coach survived only four seasons before being fired, churn that opens roughly a quarter of the league's head-coaching jobs in a typical off-season. Teams pour weeks of interviews, scheme presentations, and background checks into the choice, sometimes enlisting outside search consultants \citep{perez2025, brandt2026}, yet hires end in dismissal often enough to raise a basic question: at the moment a coach is hired, before a single game has been played, how foreseeable is the eventual outcome, whether and when the hire ends in firing, from the information the team actually has in hand?

A substantial literature studies why NFL coaches are fired, and its consistent finding is that dismissal follows poor on-field performance during a coach's tenure in the role \citep{allen2012, salaga2020, foreman2019, farnell2021}. As Section~\ref{sec:litreview} details, these studies estimate a common quantity: the firing hazard conditional on performance realized \emph{during} the tenure, measured through team-season win percentage, results relative to expectations, or player misconduct. That quantity answers a retrospective question, which in-tenure outcomes accompany a firing, and its answer, that losing leads to firing, offers a hiring team little, because it is expressed in information the team does not yet hold.

We estimate a different quantity. Fixing the information set to what is observable when the contract is signed, we ask how much of an eventual involuntary separation is foreseeable ex ante. This is not the in-tenure question asked with earlier data; it is a different estimand. The distinction is not cosmetic, because the two can diverge: a record that strongly accompanies firing once it has been realized may be close to unforecastable beforehand. It is this second, decision-relevant magnitude that bears on how teams should evaluate candidates and on whether the league's heavy coaching turnover reflects correctable misjudgment or largely irreducible uncertainty \citep{bryson2024}.

This framing yields two contributions. First, we bound how much of an eventual firing is foreseeable at the time of hire, and find the ceiling low: most of what determines a coach's fate is realized only after he is hired. Second, we show that what foreseeable signal exists is structural rather than statistical, residing not in the detailed on-field record that hiring searches scrutinize but in how the hire is made, most clearly in whether the coach is promoted from within. We develop this internal-promotion hazard as a question about the coaching labor market, reading it through the economics of internal versus external hiring and testing the leading alternative explanation, that struggling franchises both promote from within and part with coaches quickly, against observable measures of franchise distress.

\section{Literature Review}
\label{sec:litreview}

\subsection{Managerial dismissal in the NFL coaching labor market}

The dismissal of a head coach is a personnel decision, and NFL coaching turnover has accordingly drawn sustained attention in the economics of professional-sport labor markets. The consistent message of this work is that dismissal is driven by performance. \citet{allen2012} estimate probit models of turnover over 1979--2006 and exploit the early-1990s introduction of the salary cap and free agency as a natural experiment, finding that the probability of dismissal rose in the post-cap era and that measures of player talent, coaching talent, and on-field performance gained importance as determinants of turnover, indicative of heightened expectations of head coaches. \citet{salaga2020} apply survival methods to 208 head coach employment spells from 1985 to 2018 and conclude that dismissal is driven largely by raw, and to a lesser degree relative, in-season performance; their fixed covariate of central interest is the coach's race, and they do not ask how much of a separation is foreseeable from the pre-hire record. \citet{foreman2019} fit a hazard model to 505 team-seasons (2000--2015) and show that player deviance, measured through fines, penalties, and off-field legal incidents, raises the dismissal hazard. \citet{farnell2021} examine the full coaching hierarchy and find that young, experienced, well-performing coordinators tend to be promoted, while older and poorly performing coaches tend to be fired, with race playing little role in either decision.

What unites these studies is an estimand: each models the firing hazard conditional on performance realized \emph{during} the tenure. They therefore characterize which in-season outcomes accompany a firing, not what, if anything, is foreseeable about it at the moment of hire, and none conditions the outcome on only the information available when the contract is signed. Our question is instead one of forecasting from pre-realization information, in the manner of \citet{wolfson2011}, who predict NFL quarterback success from pre-draft observables and find much of the variation in later performance unforeseeable in advance.

\subsection{Internal versus external hiring}

Our second result speaks to a different margin: the \emph{source} of a hire rather than its in-season record. Whether to fill a senior post by promoting from within or by recruiting from outside is a classic question in the personnel economics of organizations. \citet{chan1996} shows that the contest between an inside and an outside candidate is optimally tilted toward the insider, so that an outsider is recruited only when distinctly superior; external hires therefore clear a higher ability bar than internal promotions, a selection difference with direct implications for who survives. In the NFL specifically, \citet{farnell2021} document that teams fill their coaching positions through a largely internal promotion ladder, and \citet{roach2016} finds that prior head coaching experience, a salient pre-hire credential, does not improve subsequent team performance. Whether coaches promoted from within go on to survive longer or shorter than outside hires is not settled, and our hazard-ratio estimates (Section~\ref{sec:hazard}) and the selection probe of Section~\ref{sec:probe_results} take it up directly.

\subsection{Survival modeling of managerial tenure and the case for Cox}
\label{sec:litcox}

A coaching spell is naturally a time-to-event outcome: it has a start, an end that is either observed (a departure) or not yet observed (an ongoing tenure), and a duration whose distribution is the quantity of interest. Survival analysis is the standard framework for such data for several reasons. It models the full time-to-event distribution, the hazard of firing at each point in a tenure, rather than collapsing the outcome to whether a coach was fired or to an average tenure length; it accommodates right-censoring, the coaches still employed at the end of the study window, so those incomplete spells inform the fit instead of being discarded as a regression on observed tenure would require; and it expresses covariate effects as shifts in the timing of the event, not merely its occurrence.

The framework has been used across sports management. \citet{gilfix2020} study Major League Soccer manager longevity and find that performance relative to predecessors, league context, and age at hire all affect tenure duration. \citet{salaga2020} fit Cox survival models to 208 NFL coaching spells over 1985--2018, but the on-field performance they use is realized during the tenure rather than fixed at hire.

Within survival analysis, the Cox proportional-hazards model \citep{cox1972} is among the most widely used time-to-event models across the empirical sciences, and we adopt it as our primary model. It is semi-parametric, leaving the baseline hazard (the shape of firing risk over the course of a tenure) unspecified, which avoids committing in advance to a parametric clock for when coaches are fired; its coefficients exponentiate into interpretable hazard ratios, which suits a study whose substantive interest is in \emph{which} factors foreshadow firing; and it is the established choice in the cognate sports-management literature \citep{salaga2020, gilfix2020}, keeping our results comparable to prior work. Whether this interpretability costs any accuracy is an empirical question we address in Section~\ref{sec:models}.

\section{Methods}

We study how much of an NFL head coach's eventual involuntary separation is foreseeable from the information available on hiring day, using a survival model that conditions on pre-hire covariates alone. Raw coaching histories and team performance statistics are scraped from Pro-Football-Reference.com \citep{pfr} and supplemented with team and roster data covering 1970--2024. All modeling is performed in Python with \texttt{scikit-learn} \citep{pedregosa2011} and the \texttt{lifelines} survival library \citep{davidsonpilon2019}.

\subsection{Population and outcome}

The analysis covers 371 non-interim head coaching stints held by 264 distinct coaches who were hired in the modern, post-merger era of the NFL, between 1970 and 2025; 86 coaches account for more than one stint. Mid-season interim appointments that did not convert into a non-interim role are excluded, and interim appointments later retained on a non-interim basis are anchored to the season in which that tenure began. The duration of a stint is the number of seasons the coach held the position; longer tenures correspond to better regular-season records (Appendix~\ref{app:winrate}).

\subsubsection{Competing-risks event coding}
\label{sec:coding}

We distinguish three terminal states for each stint, summarized in Table~\ref{tab:events}. The event of interest is \emph{involuntary non-retention}: a firing, a forced resignation, a face-saving ``mutual'' parting, or a performance-driven non-renewal. The two complementary states are a \emph{voluntary-while-viable} departure, which we treat as a competing event, and a tenure still ongoing at the data boundary, which is administratively censored.

Classifying departures is difficult because the labels organizations attach to a coaching change are unreliable: a fired coach is routinely described as having ``stepped down,'' and a ``mutual parting'' can be a dismissal in all but name. We therefore default to firing, coding a stint as involuntary non-retention unless we found positive, documented evidence that the coach left on his own terms while still employable. A departure is coded as voluntary, and hence censored, only when contemporaneous reporting placed it in one of three pre-specified categories:
\begin{quote}
\begin{enumerate}
\item[(A)] the coach left to take another head coaching position while still wanted by the incumbent team;
\item[(B)] the coach retired while still successful and retained (a ``retirement on top''); or
\item[(C)] the coach departed for documented health, family, or other non-football reasons.
\end{enumerate}
\end{quote}
A losing record in the coach's final season is itself disqualifying: such a stint is coded as a firing however the parting was described, and every one of the 27 voluntary departures we credit had a final-season record at or above .500. Each was adjudicated individually, with a one-line justification, a confidence level, and a source recorded in Appendix~\ref{app:nonfiring}. The remaining 32 non-firings are coaches still active at the 2025 data boundary.

\begin{table}[!ht]
\caption{Competing-risks coding of the 371 head coaching stints. The firing hazard treats voluntary departures and active coaches as censored; the cumulative incidence function (Section~\ref{sec:cif}) treats voluntary departure as a competing event.}
\label{tab:events}
\small
\begin{tabular}{llcc}
\toprule
Code & Terminal state & Count & Share \\
\midrule
1 (event) & Involuntary non-retention (firing) & 312 & 84.1\% \\
2 (competing) & Voluntary-while-viable departure & 27 & 7.3\% \\
0 (censored) & Active at data boundary & 32 & 8.6\% \\
\midrule
& Total & 371 & 100\% \\
\bottomrule
\end{tabular}
\end{table}

The two quantities we estimate call for voluntary exits to be handled in two different ways, which is standard in competing-risks analysis \citep{austin2016}. For absolute risk, a voluntary departure is a competing event: such a stint can never subsequently end in a firing, so censoring it na\"ively would overstate the cumulative firing incidence (Section~\ref{sec:cif}). For the cause-specific hazard, which asks what raises the rate of firing among coaches still at risk of it, the correct handling is the reverse, with voluntary departures and active coaches both censored. Each quantity is therefore estimated under the censoring convention appropriate to it.

\subsection{Pre-hire features}

The analysis draws on 74 candidate pre-hire features for each hire, all fixed at the time of hire and organized into eight families on two tiers: the coach being hired and the situation he inherits (Table~\ref{tab:feature_categories}). From these candidates the model selects a smaller set by stability selection (Section~\ref{sec:selection}). The description that follows concerns this candidate set.

\begin{table}[!ht]
\caption{The 74 candidate pre-hire features on two tiers: attributes of the coach being hired, and attributes of the situation he inherits. ``\#'' is the number of features each family contributes to the candidate set. All values are measured at or before the hire date.}
\label{tab:feature_categories}
\small
\begin{tabular}{@{}l c l@{}}
\toprule
Family & \# & Description \\
\midrule
\multicolumn{3}{@{}l}{\emph{The coach (who is hired)}}\\
Core experience & 8 & Age at hire, prior HC stints, years by coaching level and role \\
Career structure & 16 & Employer history, role and side exposure, internal-hire indicator \\
Prior unit performance & 26 & Pooled, orientation-signed unit statistics from prior roles \\
Recent form & 5 & Percentile and trajectory of the coach's recent unit and team performance \\
\addlinespace
\multicolumn{3}{@{}l}{\emph{The situation (what is inherited)}}\\
Hiring-team results & 3 & Turnovers and playoff appearances over the prior two seasons \\
Inherited team quality & 4 & Team rating (SRS) components and schedule strength \\
Inherited roster & 10 & Roster and starter value, age, experience, retention, draft capital \\
Organizational instability & 2 & Predecessor tenure, distinct HCs in the prior decade \\
\midrule
Total & 74 & \\
\bottomrule
\end{tabular}
\end{table}

Several inherited-quality features use the Simple Rating System (SRS), reported by Pro-Football-Reference \citep{pfr}, which rates a team by its average point margin adjusted for strength of schedule and expresses it in points per game relative to a league-average team. Its offensive and defensive components divide that rating into points generated and points prevented.

\subsubsection{Pooled unit representation}

A coach's prior on-field production comes from several roles, as an offensive coordinator, a defensive coordinator, or a head coach, each speaking to a different side of the ball. Rather than carry a separate, mostly empty statistical block for every role and side, we pool a coach's prior role-seasons into a single oriented \emph{unit} block of 33 statistics, of which 26 enter the candidate feature set after collapsing redundant pairs. Each season is mapped to the side of the ball the role controlled and signed so that larger is always better: an offensive coordinator's season contributes his offense, a defensive coordinator's the suppression of the opposing offense, a head coach's both, with turnovers and penalties sign-flipped. Per-season values are standardized against that season's league distribution,
\begin{equation}
z_{i,s,t} = \frac{x_{i,s,t} - \mu_{s,t}}{\sigma_{s,t}},
\label{eq:zscore}
\end{equation}
where $x_{i,s,t}$ is statistic $s$ for coach $i$'s relevant unit in season $t$ and $\mu_{s,t}, \sigma_{s,t}$ are that season's league mean and standard deviation, and are then pooled across the coach's qualifying role-seasons. This makes a coordinator who ran a top-decile unit in 1985 comparable to one who did so in 2020 and lets a coach's full body of work, rather than one role in isolation, inform the estimate.

\subsubsection{Missing data}

The feature design already removes one kind of missingness. Because the pooled unit block collapses all of a coach's roles into one oriented block, a coach who never held a given role produces no empty cells; and binary career facts, such as whether a coach ever served as an NFL coordinator on a given side, are encoded directly as indicators, so that ``never held the role'' is recorded as information rather than left as a value to guess. The missingness that remains is concentrated in the advanced efficiency and drive statistics the league only began recording in the 1990s, and so falls almost entirely on coaches whose qualifying seasons predate those measures. We impute it by truncated singular value decomposition, approximating the feature matrix by a low-rank factorization $\mathbf{X} \approx \mathbf{U}\boldsymbol{\Sigma}\mathbf{V}^\top$ re-fit within each resample on training data alone (Section~\ref{sec:eval}).

\subsection{Survival models}
\label{sec:models}

Our primary model is the Cox proportional-hazards model \citep{cox1972}, which specifies the firing hazard for stint $i$ as
\begin{equation}
\lambda_i(t \mid \mathbf{x}_i) = \lambda_0(t)\,\exp(\mathbf{x}_i^\top \boldsymbol{\beta}),
\label{eq:cox}
\end{equation}
leaving the baseline hazard $\lambda_0(t)$ unspecified. For predictive evaluation the model is fit with elastic-net regularization, whose penalty strength and $\ell_1$ mixing are chosen by coach-level cross-validated concordance over a grid (penalizer $0.5$ and $\ell_1$ ratio $0.25$ at the optimum; full search spaces in the online supplement). For the reported inference (Section~\ref{sec:hazard}) we refit without penalization, as detailed there.

To verify that the proportional-hazards form does not sacrifice predictive performance, we benchmark Cox against a family of accelerated-failure-time (AFT) models (Weibull, log-normal, and log-logistic) and against a gradient-boosted survival model, XGBoost with the \texttt{survival:aft} objective \citep{chen2016}. The AFT and boosted models serve only as a robustness check on the choice of primary model; all interpretation is drawn from the Cox fit.

\subsection{Cox-native stability selection}
\label{sec:selection}

From the 74 candidate features we select a robust subset by stability selection \citep{meinshausen2010} driven by a Cox-native importance criterion, so that selection, prediction, and interpretation all rest on the same model class and there is no cross-model importance disagreement to adjudicate. On each of $B = 2000$ subsamples drawn at the \emph{coach} level (half the coaches, sampled without replacement so no coach is split across the partition), we re-impute the features, fit a ridge Cox model, and rank features by the magnitude of their standardized coefficients. Writing $\hat{S}^{(b)}$ for the $q = 10$ highest-ranked features on subsample $b$, the selection frequency of feature $k$ is
\begin{equation}
\hat{\Pi}_k = \frac{1}{B}\sum_{b=1}^{B} \mathbf{1}\!\left\{k \in \hat{S}^{(b)}\right\},
\label{eq:stability}
\end{equation}
and we treat a feature as robustly selected when $\hat{\Pi}_k \geq \pi$ at threshold $\pi = 0.6$. Because a single-covariate model would leave the leading hazard ratio exposed to confounding by correlated pre-hire factors and would yield no adjusted estimate for the next-strongest effects, the reported hazard and prediction models adjust the leading effect for the features just below the threshold rather than fitting it alone; how many features that involves is determined by the selection results (Section~\ref{sec:selection_results}), and we report each feature's selection frequency alongside its estimate.

We set the retention threshold at $\pi = 0.6$, the lower end of the $[0.6, 0.9]$ range \citet{meinshausen2010} recommend, which together with a small shortlist length $q$ controls the expected number of falsely selected features; the formal error bound and the sensitivity of the selection profile to $q$ are reported in the online supplement. As an external check on the stable set, we also rank the 74 features by their Cox importance and choose how many to retain by maximizing coach-level cross-validated concordance, and Section~\ref{sec:selection_results} compares that cross-validated set against the stable set.

\subsection{Evaluation}
\label{sec:eval}

We evaluate the models by repeated leakage-free resampling. All predictive metrics are averaged over 50 independent train/test splits. Each split is performed at the coach level: every stint belonging to a given coach is assigned entirely to either training or test, so that a coach hired more than once (for example, a coordinator later rehired elsewhere) cannot appear on both sides. Within each split, SVD imputation is fit on the training partition and applied to the test partition, feature selection is recomputed, and the model is trained, so that no step has access to test-set information.

\subsubsection{Discrimination and calibration}

Survival models are conventionally judged on two distinct properties. Discrimination is whether the model ranks subjects in the correct order of risk; calibration is whether its predicted probabilities are numerically accurate. We report both.

For discrimination we use the concordance index, the survival-analysis generalization of the area under the ROC curve (AUROC). It is the proportion of comparable stint pairs whose predicted risk ordering matches their observed firing order, and so ranges from 0.5, indicating discrimination no better than chance, to 1.0, indicating that the coach fired sooner is always assigned the higher risk; a value of 0.60 indicates that the model correctly orders the firing times of two randomly drawn coaches 60\% of the time. We report two variants. Harrell's $C$ \citep{harrell1996} is the standard estimator. Because it can be biased upward under heavy censoring, we also report Uno's $C$ \citep{uno2011}, which reweights pairs by the inverse probability of remaining uncensored to remove that bias.

Discrimination concerns only the ordering of predictions, not whether the predicted probabilities are themselves correct. For calibration we use the time-dependent Brier score with inverse-probability-of-censoring weights \citep{graf1999}. At each horizon it is the mean squared difference between the predicted survival probability and the observed survival status, so that lower values are better, with zero denoting perfect prediction and 0.25 the score of an uninformative model that assigns every coach a 50\% survival probability. Averaging this score over the observable horizon yields the integrated Brier score, a single summary of probabilistic accuracy that is the survival analog of the Brier score for binary forecasts.

\subsubsection{Significance testing}

Across the 50 resamples the naive paired $t$-test and naive standard errors are anti-conservative, because the splits share training data and are not independent. We therefore use the Nadeau--Bengio correction \citep{nadeau2003}, which inflates the resampling variance by a factor accounting for the train/test overlap (test fraction $\rho = 0.2$), and we apply it both to model comparisons and to every confidence interval for a resampled quantity, namely the concordance indices, the integrated Brier score, and the ablation discriminations; this widens those intervals by about a factor of three and a half. Confidence intervals for quantities estimated once on the full data, such as the hazard ratios of Section~\ref{sec:hazard}, are model-based and unaffected.

\subsubsection{Inference and assumptions}

The reported inference rests on a few further choices. The hazard-ratio model is fit with cluster-robust (sandwich) standard errors \citep{lin1989} grouped on coach identity, because the same coach recurs across stints and ordinary standard errors would wrongly treat those stints as independent. We test the proportional-hazards assumption, that each covariate acts with a hazard ratio which stays constant over the course of a tenure, using the Grambsch--Therneau test on scaled Schoenfeld residuals \citep{grambsch1994}; a small $p$-value would flag a hazard ratio that drifts over time. Absolute firing risk is estimated with the Aalen--Johansen cumulative incidence function \citep{aalen1978}, which gives the probability that a coach has been fired by a given season while accounting for the competing possibility of a voluntary departure.

We pair these cause-specific hazards with the cumulative incidence function rather than fitting a Fine--Gray subdistribution model \citep{finegray1999}: reporting the two together is the standard recommendation when the aim is to identify risk factors rather than to forecast incidence \citep{austin2016}.

\section{Results}

\subsection{The firing-time distribution}

Across the 371 stints, the Kaplan--Meier median firing-free tenure is four seasons (95\% CI [3, 4] seasons). A log-rank test comparing stints that began before 2000 with those that began later provides no evidence of an era difference in the firing-time distribution ($p = 0.56$), so we pool eras in all subsequent models.

\subsection{Competing-risks incidence}
\label{sec:cif}

Figure~\ref{fig:cif} shows the cumulative incidence of firing and of voluntary departure. Firing accounts for nearly all terminations at every horizon: by the end of the second season the cumulative firing incidence is 0.29 (95\% CI [0.25, 0.34]), against 0.01 ([0.00, 0.02]) for voluntary departure; by season four it is 0.62 ([0.57, 0.67]) against 0.04 ([0.02, 0.06]); and by season eight it is 0.82 ([0.77, 0.86]) against 0.06 ([0.04, 0.09]).

\begin{figure}[!ht]
\centering
\includegraphics[width=0.75\textwidth]{figures/survival_cif.png}
\caption{Aalen--Johansen cumulative incidence of firing (the event of interest) versus voluntary-while-viable departure (the competing event), with 95\% confidence bands. Firing incidence exceeds voluntary-departure incidence at every horizon.}
\label{fig:cif}
\end{figure}

\subsection{Feature selection}
\label{sec:selection_results}

Cox-native stability selection identifies a single robustly selected predictor. Only the internal-hire indicator clears the $\pi = 0.6$ retention threshold, selected in 82\% of coach-level subsamples (Table~\ref{tab:stability}, Figure~\ref{fig:stability}). The next features form a tier just below the threshold whose membership depends on the per-subsample shortlist length $q$: age at hire, prior-unit yards-per-play efficiency, years as an NFL position coach, and the hiring team's inherited defensive rating each sit between 0.44 and 0.53 at $q = 10$ and rise to between 0.67 and 0.74 at $q = 20$. The internal-hire indicator, by contrast, is selected in 82 to 92\% of subsamples across the whole range of $q$, so its standing does not depend on that choice. Following Section~\ref{sec:selection}, we carry the five highest-frequency features into the reported hazard and prediction models while reading the stable set proper as the one feature that survives the error-controlled threshold. The paired corrected resampled $t$-test (Section~\ref{sec:eval}) finds no difference in concordance between this five-feature set and the larger, cross-validated set described in Section~\ref{sec:selection} ($\Delta C = +0.002$, $p = 0.53$), so the smaller set sacrifices no discrimination.

\begin{table}[!ht]
\caption{Cox-native stability selection (2000 coach-level subsamples). Selection frequency of each leading feature at per-subsample shortlist lengths $q \in \{10, 15, 20\}$, computed from the same subsamples; the retention threshold is $\pi = 0.6$ (primary $q = 10$). Frequencies at or above the threshold are shown in bold. Only the internal-hire indicator clears the threshold across the range of $q$.}
\label{tab:stability}
\small
\begin{tabular}{lccc}
\toprule
Feature & $q = 10$ & $q = 15$ & $q = 20$ \\
\midrule
Internal hire & \textbf{0.82} & \textbf{0.89} & \textbf{0.92} \\
Age at hire & 0.53 & \textbf{0.64} & \textbf{0.74} \\
Prior-unit yards-per-play efficiency & 0.44 & 0.59 & \textbf{0.71} \\
Years as NFL position coach & 0.48 & \textbf{0.60} & \textbf{0.68} \\
Inherited defensive SRS & 0.44 & 0.57 & \textbf{0.67} \\
\bottomrule
\end{tabular}
\end{table}

\begin{figure}[!ht]
\centering
\includegraphics[width=0.85\textwidth]{figures/survival_stability.png}
\caption{Cox-native stability-selection frequencies over 2000 coach-level subsamples (per-subsample shortlist $q = 10$). Only the internal-hire indicator clears the dashed retention threshold ($\pi = 0.6$); the remaining leading features fall just below it, and the four next-highest are carried into the multivariable model alongside it.}
\label{fig:stability}
\end{figure}

\subsection{Predictive performance}
\label{sec:performance}

Table~\ref{tab:bakeoff} reports discrimination for the Cox model and the four robustness benchmarks on the stable feature set. The Cox model attains a Harrell concordance of 0.602 (95\% CI [0.560, 0.645]) and a close Uno concordance of 0.631 (95\% CI [0.511, 0.751]); the agreement of the two point estimates indicates that censoring does not materially distort the concordance. No competing model separates itself from Cox: the parametric and gradient-boosted alternatives differ from it by no more than about 0.013 in concordance, well within the resampling uncertainty. The proportional-hazards form therefore carries no predictive penalty that the resampling can detect, and we read the models as indistinguishable in discrimination at this resolution, which supports adopting the more interpretable Cox model as primary.

To locate where this discrimination comes from, we compared the full model against a covariate-free null and a series of single-block ablations under the same coach-level resampling protocol (Table~\ref{tab:ablation}). A model with no covariates scores 0.500 by construction. The entire 33-statistic unit-performance block, the bulk of a coach's prior on-field record, does not separate from chance: its concordance is 0.526 (95\% CI [0.480, 0.573]), and the inherited-situation block scores 0.494 ([0.448, 0.539]). Both intervals include 0.500, so neither block is distinguishable from chance. The two structural variables that anchor the final model are individually far more informative: the internal-hire indicator alone reaches 0.564 (95\% CI [0.532, 0.596]) and age at hire alone 0.560 (95\% CI [0.515, 0.604]), both clearing the chance floor and each recovering roughly three-fifths of the distance from chance to the full model's 0.602 (95\% CI [0.560, 0.645]). A single binary feature, whether the coach was promoted from within, thus discriminates firing better than all thirty-three prior-performance statistics combined. What predictive signal exists is thus concentrated in the circumstances of the hire and in a narrow slice of the prior record, rather than in the statistical record taken as a block.

\begin{table}[!ht]
\caption{Discrimination of single-block and null-model ablations, averaged over the same 50 leakage-free coach-level splits as Table~\ref{tab:bakeoff}. Each row reports a Cox model fit on the indicated feature set alone. The covariate-free null is 0.500 by construction; the full selected model matches the headline 0.602. Intervals use the same Nadeau--Bengio correction for resample overlap as Table~\ref{tab:bakeoff}.}
\label{tab:ablation}
\small
\begin{tabular}{lc}
\toprule
Feature set & Harrell $C$ [95\% CI] \\
\midrule
No covariates (null) & 0.500 [0.500, 0.500] \\
Coach unit performance (33 stats) & 0.526 [0.480, 0.573] \\
Inherited situation (30 features) & 0.494 [0.448, 0.539] \\
Age at hire only & 0.560 [0.515, 0.604] \\
Internal-hire indicator only & 0.564 [0.532, 0.596] \\
Full selected model & 0.602 [0.560, 0.645] \\
\bottomrule
\end{tabular}
\end{table}

\begin{table}[!ht]
\caption{Model comparison on the five-feature stable set, averaged over 50 leakage-free coach-level splits. $C$ = concordance index; 95\% CI Nadeau--Bengio corrected for resample overlap (Section~\ref{sec:eval}). Higher is better; 0.5 is chance.}
\label{tab:bakeoff}
\small
\begin{tabular}{lcc}
\toprule
Model & Harrell $C$ [95\% CI] & Uno $C$ [95\% CI] \\
\midrule
Cox proportional hazards & 0.602 [0.560, 0.645] & 0.631 [0.511, 0.751] \\
Weibull AFT & 0.615 [0.573, 0.657] & 0.643 [0.526, 0.760] \\
Log-normal AFT & 0.615 [0.576, 0.654] & 0.641 [0.525, 0.758] \\
Log-logistic AFT & 0.615 [0.575, 0.654] & 0.640 [0.524, 0.757] \\
XGBoost-AFT & 0.608 [0.567, 0.648] & 0.634 [0.515, 0.753] \\
\bottomrule
\end{tabular}
\end{table}

A concordance near 0.60 carries real but modest signal: presented with two coaches, the model orders their firing times correctly about three in five times. Calibration tells a complementary story. Figure~\ref{fig:calibration} shows the inverse-probability-weighted Brier score by horizon; the integrated Brier score is 0.176 (95\% CI [0.160, 0.191]). Against the 0.25 expected of an uninformative model that assigns every coach the same survival probability (Section~\ref{sec:eval}), the integrated score of 0.176 is a Brier skill score of about 0.30: a roughly 30\% reduction in mean squared error over a no-information forecast, meaningful but far from precise. The figure is read against the dashed 0.25 line: the vertical distance below it at each horizon is the model's probabilistic skill there. That skill is greatest in the first season, where almost no coach has yet been fired and survival is all but certain, and at long horizons, where almost every coach has been; it nearly vanishes in the third and fourth seasons, where the error peaks near 0.25 and firing decisions are most contested and least determined by pre-hire information.

\begin{figure}[!ht]
\centering
\includegraphics[width=0.7\textwidth]{figures/survival_calibration.png}
\caption{Time-dependent inverse-probability-of-censoring-weighted Brier score for the Cox model by horizon, with the integrated Brier score in the title. Lower is better; the dashed line at 0.25 marks an uninformative model, and the gap below it at each horizon is the model's probabilistic skill there. Error peaks in the contested third and fourth seasons.}
\label{fig:calibration}
\end{figure}

\subsection{Hazard ratios}
\label{sec:hazard}

Table~\ref{tab:hazard} and Figure~\ref{fig:forest} report the cause-specific firing hazard ratios from a Cox model on the five highest-frequency features with cluster-robust standard errors on coach identity. We report this model \emph{without} penalization. With five near-uncorrelated features and 312 events the unpenalized maximum-likelihood fit is stable, and unlike a penalized fit it yields the unbiased coefficients and valid confidence intervals and $p$-values that inference requires. The proportional-hazards assumption is not rejected for any covariate (Grambsch--Therneau minimum $p = 0.20$).

\begin{table}[!ht]
\caption{Cause-specific firing hazard ratios from an unpenalized Cox model on the five highest-frequency features, with cluster-robust standard errors on coach identity. HR $> 1$ indicates faster firing, HR $< 1$ slower. Continuous features are on their natural scale.}
\label{tab:hazard}
\small
\begin{tabular}{lcccc}
\toprule
Feature & HR & 95\% CI & $p$ \\
\midrule
Internal hire & 1.524 & [1.150, 2.021] & 0.003 \\
Age at hire (per year) & 1.021 & [1.003, 1.039] & 0.021 \\
Prior-unit yards-per-play efficiency & 0.739 & [0.614, 0.891] & 0.001 \\
Years as NFL position coach & 1.021 & [0.989, 1.054] & 0.198 \\
Inherited defensive SRS & 1.024 & [0.992, 1.057] & 0.147 \\
\bottomrule
\end{tabular}
\end{table}

\begin{figure}[!ht]
\centering
\includegraphics[width=0.85\textwidth]{figures/survival_forest.png}
\caption{Cause-specific firing hazard ratios (unpenalized Cox, cluster-robust 95\% confidence intervals). Significant predictors are shown in blue. The internal-hire indicator, age at hire, and prior-unit yards-per-play efficiency are distinguishable from a hazard ratio of one; the last is protective, a lower hazard for coaches whose prior units gained more yards per play.}
\label{fig:forest}
\end{figure}

Three predictors are distinguishable from a hazard ratio of one. The only one robustly selected (Section~\ref{sec:selection_results}) is the internal-hire indicator: coaches promoted from within the hiring organization face a markedly higher firing hazard than external hires (HR 1.52, 95\% CI [1.15, 2.02], $p = 0.003$). Two further effects emerge in the multivariable fit. Firing risk rises modestly with age at hire, by about 2\% per year (HR 1.021, 95\% CI [1.003, 1.039], $p = 0.021$), echoing the pattern \citet{farnell2021} report for the broader coaching hierarchy. And coaches whose prior units were more efficient in yards per play face a lower hazard (HR 0.74, 95\% CI [0.61, 0.89], $p = 0.001$), the one slice of the on-field record that carries forward into firing risk, and it does so protectively. The remaining two features, years as an NFL position coach and the inherited defensive rating, have confidence intervals spanning a hazard ratio of one and so cannot be distinguished from no effect in these data. The internal-hire indicator also captures coaches elevated from interim to non-interim status.

The broader pattern across the model is the more important result. Taken as a block, the pooled offensive and defensive unit statistics that make up the largest share of the feature set discriminate firing at a level indistinguishable from chance (Section~\ref{sec:performance}); their predictive content collapses to a single dimension, yards-per-play efficiency, once the rest of the record is competed away. The commonly cited efficiency measures, third-down conversion rate, red-zone scoring rate, total points, and takeaways among them, are nowhere among the leading selected features (Table~\ref{tab:stability}); of the entire 33-statistic block, only prior-unit yards-per-play efficiency carries forward into firing risk. What foreseeable signal exists is concentrated in how the hire is structured, above all whether the coach is promoted from within rather than brought in from outside, with a thin protective contribution from prior-unit yards-per-play efficiency, rather than in the detailed statistical record.

\subsection{The internal-hire premium under observable selection}
\label{sec:probe_results}

To assess how much of the internal-hire effect reflects the circumstances under which teams promote from within, we refit the cause-specific hazard model of Section~\ref{sec:hazard} adding five further pre-hire covariates that proxy the hiring franchise's recent distress and the quality of what the coach inherits: the number of distinct head coaches the organization employed in the prior decade, the predecessor's tenure, the inherited offensive team rating, the inherited roster's aggregate player value, and the three-season trajectory of the team rating before the hire, each already among the 74 candidates of Table~\ref{tab:feature_categories} but below the selection threshold. Table~\ref{tab:probe} reports the internal-hire hazard ratio before and after this conditioning. Adding the proxies does not attenuate the premium: the internal-hire hazard ratio is 1.59 (95\% CI [1.18, 2.13], $p = 0.002$), essentially unchanged from 1.52 in the model without them. None of the five proxies is itself distinguishable from a firing hazard of one (all $p \geq 0.33$): in these data the franchises that promote from within are not, on the observable measures, the ones whose coaches are fired sooner. The internal-hire indicator is nearly orthogonal to the proxies (variance inflation factor 1.15, against a maximum of 2.46 across the augmented design), so the premium is being adjusted for the proxies rather than redistributed among collinear terms, and the proportional-hazards assumption continues to hold for it (Grambsch--Therneau $p = 0.11$). The premium is also not an artifact of the sample's left boundary: dropping the nine coaches who were sitting head coaches in 1969 and continued through the 1970 AFL--NFL merger, whose earlier histories the data cannot see, leaves it essentially unchanged (HR 1.55, 95\% CI [1.16, 2.06], $p = 0.003$).

Conditioning on observable distress therefore does not explain the internal-hire premium away. What remains is the possibility of selection on factors we do not observe, which we take up in Section~\ref{sec:discussion}.

\begin{table}[!ht]
\caption{The internal-hire firing hazard ratio before and after conditioning on five observable pre-hire proxies for franchise distress and inherited quality. Both are unpenalized cause-specific Cox fits with cluster-robust standard errors on coach identity (Section~\ref{sec:hazard}); Model 2 adds the five proxies to the five-feature stable set. The premium persists; none of the added proxies is distinguishable from a hazard of one.}
\label{tab:probe}
\small
\begin{tabular}{lccc}
\toprule
 & HR & 95\% CI & $p$ \\
\midrule
\multicolumn{4}{@{}l}{\emph{Internal-hire indicator}}\\
Model 1: five-feature stable set & 1.524 & [1.150, 2.021] & 0.003 \\
Model 2: + five distress proxies & 1.588 & [1.184, 2.131] & 0.002 \\
\addlinespace
\multicolumn{4}{@{}l}{\emph{Added proxies (Model 2)}}\\
Distinct HCs in prior decade & 0.981 & [0.883, 1.091] & 0.729 \\
Predecessor tenure & 0.981 & [0.944, 1.019] & 0.329 \\
Inherited offensive SRS & 0.998 & [0.958, 1.040] & 0.929 \\
Inherited roster value & 0.951 & [0.738, 1.225] & 0.697 \\
SRS trajectory & 0.989 & [0.946, 1.035] & 0.638 \\
\bottomrule
\end{tabular}
\end{table}

\section{Discussion}
\label{sec:discussion}

The analysis yields a connected set of findings. The headline is a ceiling on predictability: even with the full pre-hire record in hand, only a modest part of firing is foreseeable at the time of hire. What foreseeable signal exists sits mainly in the structural circumstances of the hire, with a thin protective contribution from the yards-per-play efficiency of the coach's prior units, rather than in the broad on-field record. We develop these in turn, then draw out what they imply for hiring teams and the limits that bound the conclusions.

\subsection{Dismissal is largely unforeseeable at the time of hire}

The central finding concerns how foreseeable dismissal is in aggregate. With every pre-hire feature available, a calibrated survival model forecasts firing with a concordance near 0.60 and an integrated Brier score of 0.176 (95\% CI [0.160, 0.191]). This is genuine signal, but a modest one, and it implies that the bulk of what determines a coach's dismissal is not knowable when the contract is signed. The ceiling is a property of the outcome rather than the estimator: the benchmark AFT and gradient-boosted models agree to within sampling error, so it does not move with the choice of model. The interpretation parallels \citet{wolfson2011} on quarterback projection: much of the variation in future success is driven by factors that are simply not observable in advance.

This finding refines the prior literature. \citet{allen2012}, \citet{salaga2020}, and \citet{foreman2019} all show that performance strongly drives dismissal, but they measure performance as it is realized during the tenure, and \citet{allen2012} in particular find that what matters is performance relative to expectations rather than raw results. Our pre-hire analysis speaks to the question those in-tenure studies leave open: the performance that drives firing is largely \emph{not} anticipated by the pre-hire record; teams are not failing to act on a signal they already hold, because most of what will determine the outcome is realized only once the coach is in the role and is not encoded in anything observable on hiring day.

This points to a sharper open question. We have shown that a coach's prior on-field record does not forecast his eventual firing, but we have not tested whether his prior performance \emph{relative to expectations} forecasts his future performance relative to expectations, the quantity the dismissal literature finds decisive. A coach who consistently exceeded the expectations placed on his prior units may carry that tendency forward in a way his raw statistics do not capture. Constructing pre-hire expectation baselines and testing that mapping is a natural direction for future work.

\subsection{The internal-promotion hazard}

The clearest structural pattern in the estimates, and the only signal that survives error-controlled selection, is that coaches promoted from within the hiring organization face a markedly higher firing hazard than those brought in from outside. This is a fact about the coaching labor market, and the personnel economics of internal versus external hiring offers a reading of it. \citet{chan1996} shows that the contest for a senior post is optimally biased toward the inside candidate, so a team recruits an outsider only when he is distinctly superior to the insider available. External hires therefore clear a higher ability bar than internal promotions, and the survival gap we estimate is what that selection difference would predict: the outside hire, chosen only when clearly better, lasts longer, while the internal promotion, advanced over a lower bar, is fired sooner.

The leading alternative reading is that this describes hiring situations rather than hiring channels, since franchises often promote from within precisely when they are in trouble, and a distressed franchise might both reach inside and fire quickly. We tested the observable form of that concern directly (Section~\ref{sec:probe_results}). Conditioning the internal-hire premium on five pre-hire measures of franchise distress and inherited quality left it intact, and none of those measures forecast firing on its own, so in the part of it we can observe, the distressed-franchise explanation is empirically absent. What remains beyond reach is selection on what we do not observe, an assistant elevated under circumstances the recorded measures fail to capture. We therefore read the result as a conditional association rather than a causal effect: promotion from within carries a higher firing hazard given the observable circumstances of the hire, with the residual identification resting on unobservable selection.

The remaining foreseeable signal is weaker and points the same way, toward the identity and circumstances of the hire rather than the football. Firing risk rises modestly with age at hire, while the bulk of the on-field record and the talent of the inherited roster do not forecast dismissal at all, a pattern consistent with \citet{roach2016}, who finds that prior head coaching experience does not improve team performance. The one thread of the on-field record that carries forward is the efficiency of the coach's prior units, which is protective but the weakest and least stably selected of these signals. For a hiring team, the implication is that the detailed football credentials given the closest attention in evaluation have little power to forecast the outcome it most wants to avoid.

\subsection{Implications for hiring and the turnover debate}

These findings have a practical bearing on how teams hire. A coaching search scrutinizes a candidate's prior offenses and defenses and his statistical record closely, yet as a whole that record forecasts firing no better than chance. This is not a case for ignoring the football record, which plausibly bears on performance once a coach is in place, nor against the modest protective signal in prior-unit efficiency; it is evidence that most of the record's predictive content has been competed away by the time a candidate reaches a head-coaching interview, when the league has already sorted candidates on observable competence and little residual variation in pre-hire credentials remains to separate those who will survive from those who will not. What stays foreseeable is largely structural, above all the elevated hazard attached to internal and interim-to-non-interim promotions.

The modest ceiling also speaks to the debate over whether coaching changes accomplish what teams intend \citep{bryson2024}. That coaches exert measurable on-field effects is not in doubt \citep{cannon2025}; what our results bound is not whether coaching matters but how much of a given hire's eventual fate is foreseeable when the contract is signed. If the eventual success of a hire is mostly unwritten on hiring day, the recurring cycle of firing and replacing looks less like the correction of identifiable mistakes than a sequence of draws from a distribution teams cannot fully see in advance. The practical implication is one of calibrated expectations: much of a hire's outcome is uncertain at the time it is made, in the way our calibration results quantify, and a firing that follows need not reflect an identifiable error in the hire.

\subsection{Limitations}

A few limitations qualify these results. The hazard-ratio $p$-values are computed on the same data used to select the features, so they describe the selected model rather than serving as confirmatory tests; the structural predictors, internal hire and age, are theory-motivated and would be included on a priori grounds, which mitigates the concern for them. The internal-hire hazard ratio should likewise be read as a conditional association rather than a causal effect: it survives adjustment for observable franchise distress (Section~\ref{sec:probe_results}), but observational data cannot exclude selection on unobservable determinants of who is promoted from within. And the model uses publicly available aggregate statistics; it cannot capture interview performance, leadership, or organizational fit, the kinds of unobservable factors our ceiling result implicates.

\section{Conclusion}

Framing NFL head coach dismissal as a firing-aware, competing-risks survival problem conditioned only on pre-hire information yields a clear picture. Dismissal is only modestly foreseeable at the moment of hiring, the proportional-hazards model is as accurate as more flexible alternatives, and what foreseeable signal exists lies chiefly in the structural circumstances of the hire, above all whether the coach is promoted from within, with weaker contributions from age and from the yards-per-play efficiency of the coach's prior units, rather than in the broad football record that hiring processes emphasize. The same pre-hire, firing-aware question extends naturally to other coaching markets with frequent turnover and a public hiring record, such as Major League Soccer and the NBA \citep{gilfix2020}; whether the modest ceiling we document is particular to the NFL or general to elite coaching labor markets is a question we leave to future work. Coaching changes may be frequent, but at the moment they are set in motion, their eventual outcome is mostly still unwritten.

\section*{Data Availability}
Data were scraped from Pro-Football-Reference.com. Processed data and analysis code are available in a public online repository; the URL is withheld to preserve anonymity during peer review and will be provided on acceptance.

\bibliographystyle{apacite}
\bibliography{refs}

\clearpage

\appendix

\section{Non-Firing Event Classifications}
\label{app:nonfiring}

Table~\ref{tab:nonfiring} lists every stint coded as a voluntary-while-viable departure, the competing event in our analysis, together with the category assigned under the protocol of Section~\ref{sec:coding} and a one-line justification drawn from contemporaneous reporting. All terminated stints not listed here are coded as involuntary non-retention (firings); the conservative default-to-firing rule means a stint appears below only where positive evidence of a voluntary, while-viable exit was found. The 32 further non-firing stints not shown are coaches still active at the 2025 data boundary, who are administratively censored: Aaron Glenn, Andy Reid, Ben Johnson, Brian Callahan, Brian Daboll, Brian Schottenheimer, Dan Campbell, Dan Quinn, Dave Canales, DeMeco Ryans, Jim Harbaugh, John Harbaugh, Jonathan Gannon, Kellen Moore, Kevin O'Connell, Kevin Stefanski, Kyle Shanahan, Liam Coen, Matt LaFleur, Mike Macdonald, Mike McDaniel, Mike Tomlin, Mike Vrabel, Nick Sirianni, Pete Carroll, Raheem Morris, Sean McDermott, Sean McVay, Sean Payton, Shane Steichen, Todd Bowles, and Zac Taylor. Sean Payton appears here for his current stint and in Table~\ref{tab:nonfiring} for an earlier voluntary departure, the two being distinct stints.

{\footnotesize
\begin{longtable}{@{}l c c c p{7.6cm}@{}}
\caption{The 27 head coaching stints coded as voluntary-while-viable departures (the competing event), grouped by category. Categories follow the protocol of Section~\ref{sec:coding}: (A) departed for another head coaching position while still wanted; (B) retired while still successful and retained; (C) departed for documented health, family, or other non-football reasons. ``Final (W--L)'' is the coach's record in his final season with the team.}
\label{tab:nonfiring}\\
\toprule
Coach & Stint & Final (W--L) & Cat. & Justification \\
\midrule
\endfirsthead
\multicolumn{5}{c}{\tablename~\thetable\ (continued)}\\[2pt]
\toprule
Coach & Stint & Final (W--L) & Cat. & Justification \\
\midrule
\endhead
\midrule
\multicolumn{5}{r}{\footnotesize\itshape continued on next page}\\
\endfoot
\bottomrule
\endlastfoot
\multicolumn{5}{@{}l}{\emph{(A) Departed for another head coaching position while still wanted}}\\[2pt]
Chuck Fairbanks & 1973--1978 & 11--4 & A & Left for the University of Colorado head job \\
Jack Pardee & 1975--1977 & 9--5 & A & Resigned to take the Redskins head job \\
Bobby Ross & 1992--1996 & 8--8 & A & Left for the Lions head job; Chargers wanted to keep him \\
Bill Parcells & 1993--1996 & 11--5 & A & Left in a personnel dispute to take the Jets job \\
Mike Holmgren & 1992--1998 & 11--5 & A & Left for the Seahawks head coach and GM role \\
Al Groh & 2000 & 9--7 & A & Resigned to take the Virginia (alma mater) head job \\
Jon Gruden & 1998--2001 & 10--6 & A & Traded to Tampa Bay; remained a sought-after head coach \\
\midrule
\multicolumn{5}{@{}l}{\emph{(B) Retired while still successful and retained}}\\[2pt]
Paul Brown & 1970--1975 & 11--3 & B & Retired voluntarily on top; named his own successor \\
Sid Gillman & 1974 & 7--7 & B & Retired on top (Coach of the Year); picked his successor \\
Bill Walsh & 1979--1988 & 10--6 & B & Retired days after winning Super Bowl XXIII \\
Jerry Burns & 1986--1991 & 8--8 & B & Announced his own retirement at 64; no pressure \\
Dick Vermeil & 1997--1999 & 13--3 & B & Retired after winning Super Bowl XXXIV \\
Bill Parcells & 1997--1999 & 8--8 & B & Retired on his own terms; Belichick set to succeed \\
Jimmy Johnson & 1996--1999 & 9--7 & B & Retired (burnout); Dolphins would have retained him \\
Dick Vermeil & 2001--2005 & 10--6 & B & Retired on top, still wanted \\
Bill Cowher & 1992--2006 & 8--8 & B & Stepped down after 15 seasons; SB winner the prior year \\
Bill Parcells & 2003--2006 & 9--7 & B & Retired by his own choice; the Cowboys wanted him back \\
Tony Dungy & 2002--2008 & 12--4 & B & Retired on top after a playoff run; picked his successor \\
Bruce Arians & 2019--2021 & 13--4 & B & Stepped down as a SB winner; picked successor Bowles \\
Sean Payton & 2006--2021 & 9--8 & B & Stepped away wanting time off; team would have kept him \\
\midrule
\multicolumn{5}{@{}l}{\emph{(C) Departed for health, family, or other non-football reasons}}\\[2pt]
John Madden & 1970--1978 & 9--7 & C & Retired due to an ulcer and burnout while still wanted \\
Bill Parcells & 1983--1990 & 13--3 & C & Resigned citing health after Super Bowl XXV \\
Joe Gibbs & 1981--1992 & 9--7 & C & Retired for health and family reasons \\
Bobby Ross & 1997--2000 & 5--4 & C & Resigned mid-season citing health (blood clots) \\
Joe Gibbs & 2004--2007 & 9--7 & C & Resigned for family reasons (grandson's illness) \\
Gary Kubiak & 2015--2016 & 9--7 & C & Retired for health reasons (migraine condition) \\
Bruce Arians & 2013--2017 & 8--8 & C & Retired citing family after a 42-year career \\
\end{longtable}
}

\section{Outcome Validity: Win Rate by Tenure Length}
\label{app:winrate}

\begin{figure}[!ht]
\centering
\includegraphics[width=0.75\textwidth]{figures/threshold_winrate.png}
\caption{Mean regular-season winning percentage by tenure length, which rises from 0.31 for short stints to 0.41 for medium and 0.54 for long (Kruskal--Wallis $p < 0.001$; each adjacent step ordered by a one-sided Mann--Whitney test, $p < 0.001$). Dismissal therefore tracks on-field success and is a meaningful outcome to model.}
\label{fig:winrate}
\end{figure}

\end{document}
